\ncuAbstractPage{\ncuENtitle}{英文摘要}{ABSTRACT}
This study examines a two-stage continual ensemble setting: retaining historical (Old) models and integrating models trained on recent (New) data. Ensemble rules and feature selection are compared. Given the available evidence and experimental configuration, the study evaluates one pool expansion using fixed temporal partitions.

The case study uses public annual financial records of U.S. companies from 1999 to 2018. Each failed company's last observed year defines a retrospective proxy for a next-year event, without verified filing dates; positives comprise 0.77\% of records. Experiment 1 compares Old, New, and combined retraining for four classifiers, then evaluates static and dynamic selection using a fixed six-model XGBoost pool. Historical TabM and TabICL scores are included provisionally, pending event-target reruns. Experiment 2 adds four feature selection methods at two retention levels. Members and decision thresholds are chosen using validation data only.

Recent data benefit classifiers differently. Within the fixed pool, New3 achieves the highest mean AUC (0.8601), while the validation-selected static ensemble achieves the highest F1-score (0.1795). Both improve AUC over equal weighting of all six models after Holm correction across 165 company-cluster paired tests; the static ensemble's F1-score increase does not pass this correction. Dynamic methods do not exceed these respective mean scores, and feature selection yields no consistent improvement. Tuned single classifiers obtain descriptive scores close to or above those of the ensembles, but parameter and validation differences prevent attributing gains to ensembling. Rankings also change with the final-status label, demonstrating that conclusions depend on the target definition.

\noindent\textbf{Keywords:} continual learning, class imbalance, temporal data, ensemble learning, dynamic selection, feature selection, bankruptcy prediction
